\uline{its value on the element} @@M0@@ \uline{of} @@M1@@, \uline{is injective}; \uline{its image is the subgroup of} @@M2@@ \uline{formed by the elements} @@M3@@ \uline{such that} (\uline{with notation analogous to that of} 1.4 \uline{and} 1.5):
@@M4@@

The theorem follows immediately from Theorem 1.6 and Lemma 1.9, since @@M5@@ does not depend on the filtration chosen on @@M6@@.

The analogous theorem obtained by replacing @@M7@@ by @@M8@@ is also true (cf. 1.7)

When there is no risk of confusion, we shall denote the functor @@M9@@ by @@M10@@; this is legitimate, since @@M11@@ identifies @@M12@@ with a subcategory of @@M13@@.

\sgaStableTarget{sga6.e05.para.1-11}1.11. We denote by @@M14@@ the category of graded @@M15@@-algebras, and by @@M16@@ the functor from @@M17@@ to @@M18@@ that assigns to a graded @@M19@@-algebra @@M20@@ the @@M21@@-algebra @@M22@@ endowed with the filtration canonically defined by its grading; the ``associated graded'' functor: @@M23@@ will be denoted by @@M24@@. There is an isomorphism of functors @@M25@@.

Let @@M26@@; we denote by @@M27@@ the multiplicative group of elements of @@M28@@ with augmentation @@M29@@. Endowing it with the filtration analogous to that of @@M30@@, we thus define a functor: @@M31@@.

Furthermore, for every @@M32@@, we define a homomorphism of filtered groups: @@M33@@, by sending the series @@M34@@,
